% 1. Apresentação geral da ferramenta  
% Definir o Espeon -> Falar do seu ativo principal (guias) -> Falar de propriedades do objeto importantes -> Puxar para os relatórios
% 4. Aplicação no contexto do estudo  
%    → mostra brevemente como o ESPEON foi utilizado nesta pesquisa, mas sem detalhes técnicos.
ESPEON é um \textit{software} de monitoramento comportamental não invasivo voltado para sala de aula virtual. O produto possui baixa barreira de entrada, necessitando apenas de um navegador com suporte ao \textit{Chrome API} para sua execução, o que o torna fácil de distribuir (\textit{Chrome WebStore}) e consumir (baixo uso de recursos computacionais), sendo portanto ideal para testes em sala descomplicados e eficientes.

A ferramenta faz uso majoritariamente da \textit{API chrome.tabs} para extrair dados de navegação do discente. Basicamente, mediante concessão de acesso o \textit{software} consegue acessar o objeto \textit{tab}, que é uma dada aba do navegador. Conforme o conceito de objeto, na Programação Orientada a Objetos (POO), trata-se de uma abstração de algum elemento real, com atributos e métodos \cite{kimObjectorientedDatabasesDefinition1990}. O ESPEON faz a leitura desse objeto e processa seus valores em um \textit{backend} para construir relatórios de aula, parametrizados por tempo e estado. O tempo se transcreve em métricas de atenção, enquanto o estado está mais associado às métricas de engajamento, que são medidas pelo estado de periféricos de áudio e vídeo do aluno e o estado da atributos da aba, todos valores booleanos. 

% colocar figura com as permissões do manifest

%propriedades

%falar do armazenamento

Munido desses \textit{logs}, o \textit{backend} faz seu processamento em métricas fundamentadas em estatística descritiva e os serve via \textit{API} nas formas de \textit{Javascript Object Notation (JSON)} e \textit{PDF}. A primeira opção é particularmente versátil porque permite que docentes consumam as métricas em outras aplicações ou até mesmo em \textit{Software} Proprietário, como será abordado na seção <ainda não existe>, usando o Power BI para melhor visualizar os dados. Já as \textit{responses} em PDF foram projetadas pensando no usuário menos especialista, em um formato estático.